% Generated draft: requires mathematical and source review.
\documentclass[11pt]{article}
\usepackage[margin=25mm]{geometry}
\usepackage{fontspec}
\setmainfont{Latin Modern Roman}
\usepackage{amsmath,amssymb,mathtools,unicode-math}
\setmathfont{Latin Modern Math}
\usepackage{enumitem,needspace,xurl,hyperref}
\setlength{\parindent}{0pt}
\setlength{\parskip}{5pt}
\setlength{\emergencystretch}{3em}
\newcommand{\N}{\mathbb{N}}
\newcommand{\Z}{\mathbb{Z}}
\newcommand{\Q}{\mathbb{Q}}
\newcommand{\R}{\mathbb{R}}
\newcommand{\C}{\mathbb{C}}
\newcommand{\F}{\mathbb{F}}
\newcommand{\A}{\mathbb{A}}
\newcommand{\PP}{\mathbb{P}}
\newcommand{\E}{\mathbb{E}}
\newcommand{\eps}{\varepsilon}
\newcommand{\dd}{\,\mathrm{d}}
\begin{document}
\section*{UnsolvedMath ID: 701755}
\textbf{Source problem code:} catalogue entry
\subsection{Definitions and mathematical statement}
A \emph{\(2\)-group} means a group \(G\), possibly infinite, in which every element has finite order equal to a power of \(2\). A group is \emph{locally finite} if every finitely generated subgroup is finite. An \emph{involution} is an element \(t\) of order exactly \(2\), and its centralizer is
\[
C_G(t)=\{g\in G:gt=tg\}.
\]
\textbf{Question (N. M. Suchkov, Kourovka Notebook Problem 20.97).} Is it true that, for every \(2\)-group \(G\),
\[
\bigl(\text{for every }t\in G\text{ with }t^2=1\text{ and }t\ne1,
\ C_G(t)\text{ is locally finite}\bigr)
\quad\Longrightarrow\quad
G\text{ is locally finite}?
\]
Here the periodic meaning of \(2\)-group stated above is the conventional interpretation of the source terminology; the supplied record does not explicitly define that term.
\subsection{Short English statement}
Must a possibly infinite group whose elements all have finite orders that are powers of \(2\) be locally finite if the centralizer of each element of order \(2\) is locally finite?
\subsection{Sources}
\begin{itemize}
\item Local UnsolvedMath registry, supplied source record with permanent ID \(701755\), exact source code \texttt{catalogue entry}, titled \emph{Kourovka Notebook Problem 20.97}.
\item The supplied record attributes the problem to N. M. Suchkov, proposed in 2022, and cites \emph{The Kourovka Notebook (Unsolved Problems in Group Theory)}, 20th issue (2022), and the 21st issue edition (2026), identified as \texttt{arXiv:1401.0300v48}. The supplied collection description names E. I. Khukhro and V. D. Mazurov as editors.
\item The registry labels the problem \emph{open}; this is supplied source metadata, not an independently verified mathematical or literature assessment. The cited editions and their wording have not been independently checked here.
\end{itemize}
\end{document}
